\section{Troubleshooting}\label{app:troubleshooting}

\subsection{``Mixed inputs'' or exit code 2}

An invocation is either all program outputs (build) or all databases (merge).
Mixing the two kinds is a usage error. Run the two steps separately:

\begin{verbatim}
CJKQC *.out -o part.jsonl --fout
CJKQC part.jsonl other.jsonl -o all.jsonl --fout
\end{verbatim}

Classification is by content, not extension, so a file whose first non-blank
line begins with \texttt{\{} counts as a database.

\subsection{An input is not recognized}

A build input that matches no supported format, or is empty, aborts the run
and writes nothing. Check:

\begin{itemize}
  \item that the file is really a program output and not, for example, a
        shell script or a truncated header;
  \item that the file is not empty;
  \item that a wrapper has not hidden the program banner beyond the head
        window; try increasing \flg{--head-bytes}.
\end{itemize}

A failed or truncated \emph{job} is different: if \CJKQC{} recognizes the
format, the record is stored and reported in the summary line.

\subsection{``cut by the window seam''}

A geometry or frequency table was split between the head and tail windows.
Widen the windows and retry:

\begin{verbatim}
CJKQC huge.out --head-bytes 1048576 --tail-bytes 8388608 -o db.jsonl --fout
\end{verbatim}

The defaults cover the shipped fixture corpus, but unusually large or
unusually laid-out outputs may need more.

\subsection{Thermochemistry prints dashes}

An undefined derivation is not an error. Check for:

\begin{itemize}
  \item a missing \fld{electronic_energy};
  \item a missing geometry, or a degenerate geometry;
  \item an unknown element symbol in the geometry;
  \item a negative frequency without \flg{--makereal} or
        \flg{--ignoreimag};
  \item a non-monatomic record with no participating vibrational mode, such
        as a single-point output or a plain xyz file.
\end{itemize}

Try \flg{--makereal} to treat an imaginary mode as real, or
\flg{--ignoreimag} to drop it. Only one may be used at a time.

\subsection{``\texttt{-{}-pop} db.jsonl fails''}

\flg{--pop} takes an optional level, and a space-separated value is consumed
even if it is your database path. Put \flg{--pop} after the positional inputs,
or use the equals form:

\begin{verbatim}
CJKQC db.jsonl --pop
CJKQC --pop=g db.jsonl
\end{verbatim}

The same value-eating behavior applies to \flg{--extrapolate}; place it after
the positional inputs.

\subsection{A correction does not attach}

Matching is exact and case-sensitive on \fld{name}. Common causes:

\begin{itemize}
  \item the correction record's name differs in case or spelling;
  \item the run record has no \fld{name};
  \item the correction record is not \fld{ok}, or lacks
        \fld{electronic_energy}, so it is not a candidate;
  \item the correction record's name is absent.
\end{itemize}

Read the standard-error warnings: skipped candidates and unmatched candidates
are reported individually or as an aggregate count.

\subsection{Extrapolation refuses to run}

Extrapolation rejects ambiguity rather than producing a silent mix:

\begin{itemize}
  \item two candidates with the same name in one source database;
  \item contributors of one name with different programs or versions;
  \item contributors of one name with exactly equal electronic energies,
        which usually means the same file entered two sources or the same
        path was passed twice;
  \item a run record whose name is missing from at least one source. A
        partially covered name is warned and not corrected.
\end{itemize}

Every source path must be a database. Use \flg{--add-correction} for a single
ordinary correction; a run cannot use both.

\subsection{``exists; not interactive; use \texttt{-{}-fout} to force''}

An output file already exists and standard input is not a terminal, so
\CJKQC{} will not ask. Either run interactively and answer \texttt{y}, or add
\flg{--fout}. The same rule applies to \texttt{movie.xyz}. A refusal at
\flg{--out} ends the run before \texttt{movie.xyz} is considered.

\subsection{A formation block or solvation block is missing}

Formation and solvation need monomers in the run's final record set:

\begin{itemize}
  \item a monomer is an \fld{ok} record with exactly one component at count 1;
  \item if a filter removed the monomer, it cannot serve as a reference;
  \item for solvation, the solvent code must appear in at least one record of
        the section's record set;
  \item a state whose solvent monomer is absent or has no level value prints a
        dash.
\end{itemize}

For non-water solvents, supply \flg{--psolvent}; the water saturation
correlation is not applied to other solvents.

\subsection{Sorting by name does not work}

Use \flg{--sort filename}. The accepted value is \fld{filename}, which sorts
by the record's \fld{name}. The help text's phrase ``name'' refers to that
name, not to a separate spelling.

\subsection{Populations or formation values changed after filtering}

Populations, formation references, and solvation distributions are computed
over exactly the records the run prints. Adding a filter changes the set, so
the normalization and the selected reference monomers can change too. This is
intended: these are ensemble properties of the printed set, not stored
per-record values.
